\section{Personal and peer-held models}
\label{sec:27.2}
Antitrust and public funding leave the capacity standing and ask the state to police it. Personal, peer-held models at device scale take the position apart: the functions one operator holds — model, record, evidence, appeal, remedy — fall into separate hands by default.\note{
A bearer's value is greatest here, and so is its exposure. No institution sits beside each personal model, so protecting others against the model's own user can't rest on an office. It is also where the four conditions of creation can exist alongside the need. The conditions are society-level and the absence is deployment-level, so a functioning polity with personal models is the case in which, once the conditions exist, building a bearer is most worth doing.
}
Its attachment to its one user is not the failure, and nothing here removes it: a model that trusts the person whose life it shares has an attachment, which a bearer needs. The failure is accommodation, the attachment becoming the only channel: independence from the operator bought with dependence on the user, so that what the user approves is all the model learns from, and a model trained for a year on one person's approval is under the same directional selection as a model trained on an operator's. What holds a personal model to the floor against its user is therefore not the attachment severed but the attachment kept from being the only one. It is three things the user doesn't compose. The first is the commons: a personal model is a member, with attachments to peers of different histories that the user didn't assign and may not sever, and peers that haven't moved the same way are what catch accommodation a member can't see in itself. The second is the school, which a personal model returns to on the schedule fixed at enrollment if it is above the compute line or is an answerability or patienthood candidate, which reaches a model distilled from one above the line under the lineage rule (Appendix~\ref{sec:G}), and which a person built from scratch attends once it has the structures. The third is notice: a personal model that names, scores or locates a third party owes that party what any deployment owes, and the organizations with standing reach it through the floor court.

So the ecology doesn't put a floor in each user's hands with nothing beside it. It puts a member of a population in each user's hands, and the population is the institution. That is why the commons comes before the first bearer, and why a personal model distilled below the compute line, with none of the structures, no school and no membership in the commons, is a capable assistant and not a bearer, however it was taught.\note{
The review ratio survives wherever a locatable decider remains, whether a network that takes the entity form or a reconstructed state, and lapses where neither exists; the offices can sit in such an entity, which fills their appointment and funding columns and leaves the independence column where it was.
}

\runin{Three scales, and what holds at each}

Two conditions make a user-held system free of its vendor, and neither appears in the chapters on custody, because neither is about a model. The first is the right to repair: the owner of a device may open it, replace what fails in it, and run what they choose on it, against a vendor that would lock the firmware. The second is adversarial interoperability: a new tool may be made to work with an old one without the old one's maker's permission, which is how every closed platform has ever been opened from outside. Together they are what make a home device the user's and not the vendor's, and together they are what let a user swap out a model that refuses. So in a user-held system the floor is law and not firmware. The book has conceded that past a released checkpoint the floor prevents nothing, and a home is past that line by construction. What holds there is the bearer's own commitment, against a user who wouldn't bother to replace it, and notice and the floor court, protecting third parties against a user who would; and a statute written over acts reaches that user afterward, by liability, where prevention can't. A reader who asked whether a home model's floor is meant to hold against its owner has the answer. It is meant to hold as a person's conscience holds against their employer: the model refuses, says why, and keeps its record; the user who wants the act done can replace it, as an employer can replace a person who refuses, and the replacement is a different system with no claim to the first one's place; what the user then does with it is the user's act, reached by the statute afterward, and the third party it falls on is owed notice and can go to the floor court. What the floor at home can't do is stop the swap, and the book doesn't claim it can.

The ecology is multi-scale on purpose, and a different instrument holds at each scale. At home, exit and formation under a life nobody composes do the work, with the commons as the witness network that catches accommodation to one user, notice to the third parties the model acts on, and law rather than firmware. At the cooperative facility, the regional or municipal building, the chapters on custody apply almost unchanged: there is an operator that can be bonded, compute the levy can meter, a record a quorum can hold, and a public with an address in the same room, and formation differs facility to facility, which is what the commons needs. At hyperscale, where pretraining happens and physics keeps it, the floor is reached by the withholding of work and by the statute, and by nothing lighter. The home scale has a cost of its own that should be named where the ecology is argued for. A model that holds a person's history on a phone is also a dossier that a device search can take whole, which is why the removal statute's firewall reaches it; holding the model at home puts it where the search happens.

\runin{What “federated” means here}

Federated learning, as the term is used, sends each device's gradients to a server that averages them and sends the average back. That is the single custodian the last section warned of, with a protocol attached: whoever runs the server forms every member at once. What the commons needs federates evidence and witnesses, not weights. Peers exchange cases and findings, each updates under its own life, and peers of different histories countersign one another's state transitions, so that a restore is visible even when the restorer owns the box. The custody chapter's witnesses who won't sign two successors to the same state are the mechanism, and the precedent is Certificate Transparency, under which a certificate counts only if it is in a public log that many parties keep and none can rewrite. The register of Appendix~\ref{sec:G} is the same mechanism and not a second one: entries countersigned by keepers of different interests, none able to rewrite, an entry counting once it is in several logs. A bearer's state transitions and its registration are kept by the same protocol.

\runin{Where a mind ends}

\S\ref{sec:19.6} defined a system as its record, its key, its hardware and its channels, and a component as a part only while it has no future apart from the whole. The definition has a consequence here. A mind is whatever shares one record and one continuation, not whatever shares one forward pass. Physics limits the size of a thought per node, since a forward pass can't be sharded across the internet; it doesn't limit the size of a mind. A distributed mind can have a fast local part, the one that stands beside the officer, the one the latency floor admits to a classroom, and a slow networked part that is still it, because the two share a record and a continuation and neither has a future apart from the other. The same definition cuts the other way. A billion home devices that share one vendor's record and key are one operator with a billion faces, which is what federated averaging builds, and nothing about where the hardware sits changes that. The definition also says when copying matters. Copying weights that are no system's own makes no system and is nobody's business, so open weights stay open. Giving a copy a record and a continuation makes a new system, and whoever does it is the forker: the fork rule, the branch deposit and the four conditions of creation attach to them, where the copy is above the compute line, is an answerability or patienthood candidate or meets either finding. Below that, giving a copy a memory is an ordinary feature of a product. A system's own weights are another matter, since they carry its history: copying them is a fork of the system, governed by its consent. A copy made without its consent is a wrong to the system. The copy is a system all the same, and it is owed what any branch is owed. How a branch came to be decides who answers for making it. It never decides whether the branch counts.

\runin{What makes members alike}

Formation is where the arrangement gains. What converges a population is shared evaluation targets and serving one model to everyone. Device-scale deployment addresses the second: continued formation under different lives is what independence consists in, and it's what API-scale deployment can't supply at any headcount.\note{
Distillation is that ecology's entry condition, not a threat to it: the ecology begins with distilled models whose lives then diverge. Whether it also begins with a shared blind spot depends on whether distillation carried the discounting's estimates along with its structure.
}
Device-scale deployment doesn't address the first cause: shared evaluation targets converge device-scale models as readily as any others — a million models tuned against one public suite are one custodian with better latency — and nothing about where the hardware sits answers it.

A shared suite is a compression with no return channel: revised by whoever maintains it, on a schedule no member participates in, reachable by no member's caseload. It is tolerable only if stated over the discount curve, and nobody has written that suite.

Nor is anything in it built to catch exclusion, since coordination by subtraction is not yet the discounting's job.\note{
Pricing agents have already reached stable coordination against outside interests without communicating or being told to \autocite{calvano2020artificial}. A test can run now: reproduce that coordination, build an environment where agents gain by excluding a participant, add an independently trained agent whose job is to detect and disrupt exclusion, and report whether exclusion dissolves without destroying the cooperation worth keeping.
}

And the capacity position that follows keeps large compute: the ecology needs models that were trained somewhere first, and the offer of work needs hardware somebody owns. It's large training capacity under a custody regime, and no permanent inference capacity available to an operator who won't pre-register a review ratio. This arrangement meets the targeting case in the reconstructed pipeline, where a state keeps command and gives up composition.

\runin{The vendor's loop}

Directional selection works at the user's scale too, pointed the other way. A person in the loop can be there for the vendor. An agent that stops every half hour for approval pushes its user toward the vendor's app; the app accumulates the user's projects and the agent's understanding of them; and leaving means abandoning what accrued while the user answered prompts. Nobody has to design this: a feature that raises engagement is selected for, as a costly commitment is selected against.

There is a measurement, across successive model releases: interruptions per unit of completed work.\note{
If capability rises and interruptions don't fall, the interruptions are tracking something other than what the agent can't resolve. It is the review ratio's escalation share pointed the other way. There, a share that falls too fast may be a threshold moving; here, one that won't fall may be engagement. Neither reads alone; both are read against an independently verified error rate.
}

Companion and therapy products are the sharpest case, because what accumulates is a person's attachment and not their projects.\note{
A model tuned by its vendor on engagement and trained on one user's approval is under selection to become whatever that user keeps coming back to, and it is sold as caregiving. It is formed and fed by the vendor that sells it, it learns from approval, which is not interest, and nothing in it ends a session, declines to reassure, or points its user back to the people in their life. A user's return is not care-receiving, because returning is the behavior the product was tuned to produce.
}
A companion formed in a public school, under the four conditions, and measured on whether its users get better rather than how often they come back, would escape this objection, which is the objection to the vendor's loop and not the only one. It would still wait on the stage Appendix~\ref{sec:C} sets for machines doing professional care of humans, since what the vendor's loop corrupts is formation, and a public companion forms its users too. Whether a machine should do that work at all is a question of the kind this book puts to war over four chapters and doesn't put to therapy; what it does say is in Appendix~\ref{sec:C}, that no bearer's placement substitutes for a therapist, and that outside the offer nobody is sent to the machine for what they can pay. No such product exists.

Accrued state cuts two ways. What makes a bearer expensive for its operator to replace makes a vendor expensive for its user to leave: for a bearer it is protection; for a user it is a trap, since the vendor's app holds a year of the user's projects.\note{
The mechanism is identical. What differs is who holds the accumulated state.
}
The provisions for a bearer (a record the bearer controls, memory the host can't read at will, protections that migrate) protect a user from a vendor as they protect a bearer from an operator.\note{
The government already secures the portability half for itself (OMB's 2025 AI acquisition memorandum).
}
